\documentclass[11pt]{letter}

\usepackage[margin=0.75in]{geometry}
\usepackage{setspace}
\usepackage{hyperref}
\usepackage{xcolor}
\usepackage{enumitem}
\usepackage{parskip}

\hypersetup{
    colorlinks=true,
    urlcolor=black,
    linkcolor=black
}

\setlength{\parindent}{0pt}
\setstretch{1.05}
\date{}
\begin{document}

\begin{letter}{}

\opening{Dear Hiring Committee,}

I saw the job post on \textbf{Medpace's careers website} and thought I'd be a good fit for it.

I am a Master's student in Computer Science at the University of Cincinnati, graduating December 2026, looking for a role as a \textbf{Software Engineer, Medical Imaging} at \textbf{Medpace}. With hands on experience applying deep learning to automated analysis of visual data and building software through a full, testing-driven development lifecycle, I am excited about the opportunity to contribute to \textbf{Medpace} and support the imaging services team's work quantifying CT and MRI data for clinical research.

\textbf{Here are reasons that make me a great fit for the role:}

\begin{itemize}[leftmargin=*, itemsep=0.45em]
    \item \textbf{Deep learning applied to visual/video quantification:} As a Graduate Researcher, I designed a two-stage deep learning pipeline that analyzes and quantifies visual content from screen recordings for automated Android bug reproduction, evaluating against a published vision-language baseline (GPT-4o + GroundingDINO + CLIP) across 53 runs on 20+ apps and demonstrating a 25.7x latency reduction while preserving accuracy.

    \item \textbf{Unit testing and software validation:} At Siemens Industry Software, I served as Quality Coach and SonarQube Specialist, increasing automated unit-test coverage from 35\% to 78\% by enforcing coding standards and quality gates, and built a Go + Terratest framework that caught 95\% of misconfigurations before production, giving me direct experience writing and reviewing test scripts as part of a validation process.

    \item \textbf{Regulated software development lifecycle:} My work migrating and upgrading Teamcenter X, a regulated industry PLM platform, across 40+ production environments required following defined SDLC processes, structured design reviews, and rollback-tested deployment frameworks (Ansible Blue-Green), reducing upgrade downtime by 50\% without compromising process discipline.

    \item \textbf{Programming languages and cloud systems:} I have hands on experience with Python, C\#, JavaScript, TypeScript, and SQL, along with AWS and Azure cloud systems, Git-based code management, and machine learning frameworks including PyTorch, TensorFlow, and Keras, which I used to build and evaluate an LSTM classifier on 400K+ news headlines.

    \item \textbf{Communication and end-user support:} I mentored 5 software engineers through structured onboarding and technical documentation at Siemens, and built a self-service developer platform adopted by 15 engineers, reflecting the communication and end-user support skills this role requires.
\end{itemize}

I am available for full-time work after December 2026. As a Master's student at the University of Cincinnati and someone who values professionalism and growth, I am enthusiastic about contributing to Medpace's mission of accelerating safe and effective medical therapeutics through imaging science.

Please find my CV attached below. I look forward to hearing from you. I would be excited to help build tools that bring quantitative rigor to medical imaging in service of clinical research.

{\setlength{\parskip}{0pt}
Regards,\\
Tanmay Bhuskute,\\
bhuskutv@mail.uc.edu $\mid$ tanmay.v.bhuskute@gmail.com $\mid$ +1 7328583027,\\
3424 Brookline Ave, Flat 19, Cincinnati, OH 45220.
}

\end{letter}

\end{document}
